\begin{table}[!htbp]
\centering
\small
\setlength{\tabcolsep}{3.5pt}
\caption{Value of $\gamma$ at which the pre-treatment steady state loses stability (Hopf bifurcation), at $\sigma_P=0.21$ and the prior centres otherwise, as a function of $\mathrm{MTT}_{\mathrm{N}}$ (from the lower to the upper limit of its search box), for the Model3 neutrophil chain ($n_{\mathrm{N}}=7$, $k_{\mathrm{N}}=4$) and the Model2 chain (3 and 2). 2.37 is the value set by the megakaryocyte lineage, which applies when the neutrophil lineage is more stable. The search box of $\gamma$ ends at 1. On a $6\times6\times6$ grid over the search boxes of $\mathrm{MTT}_{\mathrm{N}}$, $\mathrm{MTT}_{\mathrm{Pl}}$ and $\mathrm{MTT}_{\mathrm{Mo}}$ at $\gamma=1$, the largest real part of the eigenvalues is $-0.004$~d$^{-1}$.}
\label{tab:hopf}
\begin{tabular}{rrr}
\toprule
$\mathrm{MTT}_{\mathrm{N}}$ (d) & Hopf $\gamma$, Model3 chain & Hopf $\gamma$, Model2 chain \\
\midrule
0.677 & 2.37 & 2.37 \\
2 & 2.37 & 2.37 \\
3.5 & 2.37 & 2.37 \\
5 & 2.09 & 2.37 \\
8.5 & 1.48 & 2.25 \\
12 & 1.22 & 1.95 \\
15 & 1.10 & 1.82 \\
\bottomrule
\end{tabular}
\end{table}
